% ══════════════════════════════════════════════════════════════════
% Section 9 — Conclusion
% ══════════════════════════════════════════════════════════════════
\section{Conclusion}
\label{sec:conclusion}

This paper presented a systematic comparison of neural network surrogates for
hyperelastic constitutive models, covering the full pipeline from synthetic
data generation through physics-informed training to automated Fortran~90 code
emission for finite element solvers.

The main findings are:

\begin{enumerate}
  \item \textbf{Architecture comparison.}  The MLP, ICNN, and polyconvex ICNN
    architectures all achieve high accuracy ($R^2 > $
    % TODO: Insert minimum R^2 across all benchmarks.
    \textbf{[XX]}) on five hyperelastic materials (three isotropic, two
    anisotropic) when trained with the energy--stress dual loss.  The MLP
    provides the best pointwise accuracy, while the ICNN and polyconvex
    variants offer convexity and polyconvexity guarantees by construction.

  \item \textbf{Energy--stress loss.}  The dual loss combining energy
    regression with automatic-differentiation-based stress matching is
    essential for accurate gradient predictions and is the only viable
    training strategy for scalar-output architectures (ICNN, polyconvex).

  \item \textbf{Scaling behavior.}  Surrogate accuracy saturates beyond
    % TODO: Insert saturation threshold.
    \textbf{[N]} training samples for the materials considered.  A hidden
    layer width of 64 provides a good accuracy--cost balance for 3D isotropic
    inputs.

  \item \textbf{Fortran export.}  The hybrid NN--UMAT approach --- where the
    neural network evaluates the strain energy while analytical Fortran code
    computes kinematics, stress, and the consistent tangent --- produces
    self-contained subroutines compatible with Abaqus and FEAP without any
    Python runtime dependency.

  \item \textbf{FE validation.}  Single-element tests confirm that the hybrid
    NN--UMAT reproduces the analytical Cauchy stress to within
    % TODO: Insert actual error bounds.
    \textbf{[XX]\%} relative error and achieves quadratic Newton--Raphson
    convergence, validating the consistent tangent implementation.
\end{enumerate}

The \texttt{hyper-surrogate} library is released as open-source software,
providing the complete pipeline as a reusable tool for researchers and
engineers.

\subsection{Future work}

Promising directions include:
\begin{itemize}
  \item Extension to multi-fiber anisotropic models (9D+ input) and
    viscoelastic/damage formulations with internal state variables.
  \item Adaptive training data generation guided by FE simulation error
    indicators.
  \item GPU-accelerated batch evaluation of NN-based constitutive models
    across integration points.
  \item Integration with automatic model selection to choose the optimal
    architecture for a given material class.
  \item Application to experimental data fitting, where the NN surrogate
    replaces a parametric model entirely.
\end{itemize}
